\chapter{Marco conceptual}

El marco conceptual de esta tesis se organiza alrededor de cinco conceptos
articuladores: plataforma cuadrúpeda Nova Spot Micro, modelo
cinemático-dinámico, sistema de control de locomoción, aprendizaje por refuerzo
correctivo y evaluación de la caminata. Estos conceptos no se seleccionan por
afinidad temática, sino porque aparecen en el planteamiento del problema, son
necesarios para cumplir los objetivos o se usan directamente en el análisis de
resultados. En conjunto, permiten interpretar el proyecto como un proceso que va
desde una plataforma física inactiva hasta una caminata planificada,
supervisada y evaluable.

La plataforma cuadrúpeda Nova Spot Micro constituye el punto de partida material
de la investigación. El planteamiento del problema muestra que el robot
permanece inactivo por limitaciones de accionamiento, sensorización, potencia,
documentación técnica, modelado y control. En esta tesis, la plataforma no se
entiende solo como el robot disponible en el laboratorio, sino como el sistema
físico que impone restricciones a todo el desarrollo: tipo de actuadores,
necesidad de calibración, seguridad eléctrica, comunicación y capacidad de
medición. Por eso, cualquier estrategia de locomoción debe relacionarse con las
condiciones reales del Nova Spot Micro y no únicamente con resultados ideales de
simulación.

El modelo cinemático-dinámico articula la plataforma física con el diseño de
movimiento. Su inclusión se sustenta en el problema identificado: la ausencia de
un modelo matemático que describa la cinemática y dinámica del robot limita la
generación de trayectorias coherentes. La cinemática permite analizar si las
posturas, el paso y el gateo son geométricamente alcanzables por las doce
articulaciones; la dinámica permite interpretar cómo masa, inercia, gravedad,
fricción, contacto y saturación influyen en la respuesta del sistema. La
posición adoptada en esta tesis es tratar dicho modelo como una base nominal
para simulación y análisis, no como un gemelo digital identificado, porque sus
parámetros aún requieren contraste físico.

El sistema de control de locomoción conecta el modelo con la ejecución de
movimiento. En el problema se afirma que no existe una estrategia que permita
generar patrones de marcha estructurados; por ello, el control no se reduce a
enviar señales a los actuadores, sino que integra referencias de postura, paso,
gateo, límites articulares, supervisión y condiciones de seguridad. La
interpretación adoptada es conservar una marcha nominal como referencia
principal, porque permite saber qué pata debe oscilar, cuáles deben permanecer
en apoyo y qué trayectoria debe seguir cada articulación. Esta decisión hace que
la locomoción sea trazable y comparable antes de incorporar una corrección
aprendida.

El aprendizaje por refuerzo correctivo se relaciona con el proyecto porque la
pregunta problema plantea evaluar una estrategia aprendida como capa correctiva
de una marcha nominal. Este concepto amplía el sistema de control, pero no lo
sustituye. La política aprendida se interpreta como un mecanismo para proponer
ajustes pequeños sobre referencias ya planificadas, mientras el supervisor y los
límites del sistema conservan autoridad sobre la seguridad. Esta lectura es
importante para la investigación porque evita presentar el aprendizaje por
refuerzo como una solución autónoma sin restricciones y permite compararlo de
forma justa contra la marcha nominal.

La evaluación de la caminata cierra la relación entre problema, objetivos y
análisis. El proyecto no busca demostrar únicamente que el robot se desplaza,
sino establecer si la locomoción puede considerarse estable, continua y
repetible bajo condiciones controladas. Por esta razón, la evaluación se
organiza mediante tres pruebas: margen de estabilidad estática, seguimiento de
posición articular y repetibilidad del patrón de marcha. Estas pruebas conectan
el modelo, el control y el aprendizaje con resultados verificables; además,
permiten aceptar, rechazar o limitar las afirmaciones sobre el funcionamiento
del robot.

\begin{table}[H]
\centering
\caption{Conceptos articuladores del marco conceptual}
\label{tab:marco-conceptual-relacion}
\small
\setlength{\tabcolsep}{3pt}
\begin{tabular}{p{0.22\textwidth}p{0.23\textwidth}p{0.22\textwidth}p{0.23\textwidth}}
\toprule
Concepto & Presencia en el problema & Relación con los objetivos & Uso en el análisis \\
\midrule
Plataforma cuadrúpeda Nova Spot Micro & Es el sistema físico disponible, pero inactivo por limitaciones de accionamiento, sensado, potencia y documentación. & Delimita la implementación de hardware y software necesaria para ejecutar movimiento. & Condiciona alcance, calibración, seguridad, transferencia y restricciones reales. \\
Modelo cinemático-dinámico & El problema identifica la ausencia de un modelo matemático de la plataforma. & Sustenta el modelado para simulación y posterior diseño. & Permite analizar alcanzabilidad, continuidad, contacto, estabilidad y diferencias entre simulación y robot físico. \\
Sistema de control de locomoción & No existe una estrategia que genere patrones de marcha estructurados. & Integra paso, gateo, marcha nominal, aprendizaje y supervisor. & Permite estudiar si el movimiento se produce de manera coordinada y limitada. \\
Aprendizaje por refuerzo correctivo & La pregunta problema lo plantea como capa correctiva sobre una marcha nominal. & Sustenta la estrategia de control aprendida del objetivo general. & Se compara contra la marcha nominal bajo las mismas condiciones experimentales. \\
Evaluación de la caminata & El problema exige validar estabilidad y marcha, no solo producir movimiento. & Corresponde a margen estático, seguimiento articular y repetibilidad. & Define los criterios para aceptar, rechazar o limitar los resultados obtenidos. \\
\bottomrule
\end{tabular}
\end{table}
